% GENERATED FILE. Edit canonical lessons, then run npm run generate:books.
% canonical-source-hash: fnv1a64:fd438dfd090d78b0
% canonical-lessons: RU-C291-sledovatel, RU-C291-rassledovaniye, RU-C291-rassledovat, RU-C291-podozrevayemyy, RU-C291-podozrevat

\chapter{The investigation}
\label{ch:ru-the-investigation}
\hlchaptermodality{291}

\begin{chapteropening}
\textbf{By the end of this chapter:} \emph{I can talk about an investigator, an investigation, investigating, a suspect and suspecting someone.}

Complete the last lesson of chapter 291: I can talk about an investigator, an investigation, investigating, a suspect and suspecting someone.
\end{chapteropening}

\section[\texorpdfstring{\ru{следователь} (slédovatel)}{следователь (slédovatel)}]{\ru{следователь} (slédovatel) — an investigator, a detective}

\label{lesson:RU-C291-sledovatel}



\begin{quote}
Before the new one: say the Russian for to kidnap, then the Russian for to hide.
\end{quote}

\subsection*{You'll want to know: \ru{следователь}}
\textbf{\ru{следователь}} (\emph{slédovatel}) — \textquotedblleft{}an investigator, a detective\textquotedblright{}.

\begin{quote}
\ru{Следователь} \ru{долго} \ru{говорил} \ru{со} \ru{свидетелем}.
\end{quote}

\begin{quote}
\emph{Slédovatel dólgo govoríl so svidételem.}
\end{quote}

\begin{quote}
The investigator talked with the witness for a long time.
\end{quote}

Genitive \emph{\ru{следователя}}.

\begin{cousinweb}
\emph{\ru{следователь}} is built on \emph{\ru{следовать}} \textquotedblleft{}to follow\textquotedblright{}, from \emph{\ru{след}} \textquotedblleft{}a track, a trail\textquotedblright{} (EtymDB): one who follows the footsteps.
\end{cousinweb}

\begin{grammarlens}[title={What you've built}]
One more word for telling, explaining and describing at B1.
\end{grammarlens}

\subsection*{Your turn}
\begin{itemize}
\raggedright
  \item \emph{Say it:} \emph{slédovatel}
  \item \emph{Say it:} \emph{slédovatel}, once more
  \item \emph{Say it:} a sentence of your own with \emph{slédovatel}
  \item \emph{Recall:} say the Russian for to kidnap, then the Russian for to hide, then say \emph{slédovatel} again
\end{itemize}

\begin{tcolorbox}[breakable,colback=teal!4,colframe=teal!35!black,title={Before you move on}]
What does \ru{следователь} mean? (\textquotedblleft{}An investigator, a detective\textquotedblright{}.) Say it once more.
\end{tcolorbox}
\section[\texorpdfstring{\ru{расследование} (rasslédovaniye)}{расследование (rasslédovaniye)}]{\ru{расследование} (rasslédovaniye) — an investigation}

\label{lesson:RU-C291-rassledovaniye}



\begin{quote}
Before the new one: say the Russian for to hide, then the Russian for an investigator.
\end{quote}

\subsection*{You'll want to know: \ru{расследование}}
\textbf{\ru{расследование}} (\emph{rasslédovaniye}) — \textquotedblleft{}an investigation\textquotedblright{}.

\begin{quote}
\ru{Расследование} \ru{продолжалось} \ru{два} \ru{месяца}.
\end{quote}

\begin{quote}
\emph{Rasslédovaniye prodolzhálos dva mésyatsa.}
\end{quote}

\begin{quote}
The investigation went on for two months.
\end{quote}

Neuter.

\begin{cousinweb}
\emph{\ru{расследование}} is built on \emph{\ru{расследовать}} \textquotedblleft{}to investigate\textquotedblright{}, from \emph{\ru{рас}-} and \emph{\ru{след}} \textquotedblleft{}a track\textquotedblright{} (EtymDB): following the footsteps to the end.
\end{cousinweb}

\begin{grammarlens}[title={What you've built}]
One more word for telling, explaining and describing at B1.
\end{grammarlens}

\subsection*{Your turn}
\begin{itemize}
\raggedright
  \item \emph{Say it:} \emph{rasslédovaniye}
  \item \emph{Say it:} \emph{rasslédovaniye}, once more
  \item \emph{Say it:} a sentence of your own with \emph{rasslédovaniye}
  \item \emph{Recall:} say the Russian for to hide, then the Russian for an investigator, then say \emph{rasslédovaniye} again
\end{itemize}

\begin{tcolorbox}[breakable,colback=teal!4,colframe=teal!35!black,title={Before you move on}]
What does \ru{расследование} mean? (\textquotedblleft{}An investigation\textquotedblright{}.) Say it once more.
\end{tcolorbox}
\section[\texorpdfstring{\ru{расследовать} (rasslédovat)}{расследовать (rasslédovat)}]{\ru{расследовать} (rasslédovat) — to investigate}

\label{lesson:RU-C291-rassledovat}



\begin{quote}
Before the new one: say the Russian for an investigator, then the Russian for an investigation.
\end{quote}

\subsection*{You'll want to know: \ru{расследовать}}
\textbf{\ru{расследовать}} (\emph{rasslédovat}) — \textquotedblleft{}to investigate\textquotedblright{}.

\begin{quote}
\ru{Журналисты} \ru{расследовали} \ru{эту} \ru{историю}.
\end{quote}

\begin{quote}
\emph{Zhurnalísty rasslédovali étu istóriyu.}
\end{quote}

\begin{quote}
The journalists investigated this story.
\end{quote}

One form for both aspects. It takes the accusative. Present \emph{\ru{расследую}}.

\begin{cousinweb}
\emph{\ru{расследовать}} is built on \emph{\ru{рас}-} and \emph{\ru{следовать}} \textquotedblleft{}to follow\textquotedblright{}, from \emph{\ru{след}} \textquotedblleft{}a track\textquotedblright{} (EtymDB lists the verb with no further parent).
\end{cousinweb}

\begin{grammarlens}[title={What you've built}]
One more word for telling, explaining and describing at B1.
\end{grammarlens}

\subsection*{Your turn}
\begin{itemize}
\raggedright
  \item \emph{Say it:} \emph{rasslédovat}
  \item \emph{Say it:} \emph{rasslédovat}, once more
  \item \emph{Say it:} a sentence of your own with \emph{rasslédovat}
  \item \emph{Recall:} say the Russian for an investigator, then the Russian for an investigation, then say \emph{rasslédovat} again
\end{itemize}

\begin{tcolorbox}[breakable,colback=teal!4,colframe=teal!35!black,title={Before you move on}]
What does \ru{расследовать} mean? (\textquotedblleft{}To investigate\textquotedblright{}.) Say it once more.
\end{tcolorbox}
\section[\texorpdfstring{\ru{подозреваемый} (podozreváyemyy)}{подозреваемый (podozreváyemyy)}]{\ru{подозреваемый} (podozreváyemyy) — a suspect}

\label{lesson:RU-C291-podozrevayemyy}



\begin{quote}
Before the new one: say the Russian for an investigation, then the Russian for to investigate.
\end{quote}

\subsection*{You'll want to know: \ru{подозреваемый}}
\textbf{\ru{подозреваемый}} (\emph{podozreváyemyy}) — \textquotedblleft{}a suspect\textquotedblright{}.

\begin{quote}
\ru{Подозреваемый} \ru{молчал}.
\end{quote}

\begin{quote}
\emph{Podozreváyemyy molchál.}
\end{quote}

\begin{quote}
The suspect said nothing.
\end{quote}

A participle that declines like an adjective; a woman is \emph{\ru{подозреваемая}}.

\begin{cousinweb}
\emph{\ru{подозреваемый}} is the passive participle of \emph{\ru{подозревать}} \textquotedblleft{}to suspect\textquotedblright{}, built on \emph{\ru{под}-} and the old root \emph{\ru{зр}-} \textquotedblleft{}to look\textquotedblright{}, as in \emph{\ru{зрение}} \textquotedblleft{}sight\textquotedblright{}, Old East Slavic \emph{zĭrěti} \textquotedblleft{}to see\textquotedblright{} (EtymDB): the one watched from below.
\end{cousinweb}

\begin{grammarlens}[title={What you've built}]
One more word for telling, explaining and describing at B1.
\end{grammarlens}

\subsection*{Your turn}
\begin{itemize}
\raggedright
  \item \emph{Say it:} \emph{podozreváyemyy}
  \item \emph{Say it:} \emph{podozreváyemyy}, once more
  \item \emph{Say it:} a sentence of your own with \emph{podozreváyemyy}
  \item \emph{Recall:} say the Russian for an investigation, then the Russian for to investigate, then say \emph{podozreváyemyy} again
\end{itemize}

\begin{tcolorbox}[breakable,colback=teal!4,colframe=teal!35!black,title={Before you move on}]
What does \ru{подозреваемый} mean? (\textquotedblleft{}A suspect\textquotedblright{}.) Say it once more.
\end{tcolorbox}
\section[\texorpdfstring{\ru{подозревать} (podozrevát)}{подозревать (podozrevát)}]{\ru{подозревать} (podozrevát) — to suspect}

\label{lesson:RU-C291-podozrevat}



\begin{quote}
Before the new one: say the Russian for to investigate, then the Russian for a suspect.
\end{quote}

\subsection*{You'll want to know: \ru{подозревать}}
\textbf{\ru{подозревать}} (\emph{podozrevát}) — \textquotedblleft{}to suspect\textquotedblright{}.

\begin{quote}
\ru{Я} \ru{подозреваю}, \ru{что} \ru{он} \ru{нас} \ru{обманул}.
\end{quote}

\begin{quote}
\emph{Ya podozreváyu, shto on nas obmanúl.}
\end{quote}

\begin{quote}
I suspect that he deceived us.
\end{quote}

An imperfective verb. \emph{\ru{Подозревать} \ru{кого}-\ru{то} \ru{в}} + prepositional: \emph{\ru{его} \ru{подозревают} \ru{в} \ru{краже}}.

\begin{cousinweb}
\emph{\ru{подозревать}} is built on \emph{\ru{под}-} \textquotedblleft{}under\textquotedblright{} and the old root \emph{\ru{зр}-} \textquotedblleft{}to look\textquotedblright{}, as in \emph{\ru{зрение}} \textquotedblleft{}sight\textquotedblright{}, inherited through Old East Slavic \emph{zĭrěti} \textquotedblleft{}to see\textquotedblright{} (EtymDB): to watch from below. EtymDB lists the verb with no parent.
\end{cousinweb}

\begin{grammarlens}[title={What you've built}]
One more word for telling, explaining and describing at B1.
\end{grammarlens}

\subsection*{Your turn}
\begin{itemize}
\raggedright
  \item \emph{Say it:} \emph{podozrevát}
  \item \emph{Say it:} \emph{podozrevát}, once more
  \item \emph{Say it:} a sentence of your own with \emph{podozrevát}
  \item \emph{Recall:} say the Russian for to investigate, then the Russian for a suspect, then say \emph{podozrevát} again
\end{itemize}

\begin{tcolorbox}[breakable,colback=teal!4,colframe=teal!35!black,title={Before you move on}]
What does \ru{подозревать} mean? (\textquotedblleft{}To suspect\textquotedblright{}.) Say it once more.
\end{tcolorbox}
